\section{Conclusion}

We studied a restricted class of total completions of a partial algebra: the
original carrier is preserved pointwise and only finitely many new elements may
be added.  The main theorem shows that this class discriminates every finite
subset of the free compatible completion by a single homomorphism.  The proof
is direct and uses only the finite set of subterms occurring in the elements to
be distinguished.

When the base is infinite, this statement is not obtained from pairwise
separation by the usual finite-product argument, because finite complement over
the diagonal copy of the base is not preserved by products.  This explains why
the direct finite-subterm construction is the relevant one in the present
setting.

The same family of finite-complement completions defines a separated uniformity
on the free completion.  For finite bases its Cauchy completion is proper
exactly when the original algebra is genuinely partial.  For arbitrary bases,
a context fixing every base element gives a simple sufficient condition for
properness.

The intended contribution is modest: it is the finite-complement restriction
and the resulting discrimination and completion statements.  The surrounding
free-completion, residual, finite-state and Cauchy machinery is standard and has
been presented as such throughout.
